\documentclass[legal]{polinetwork}

\pnTitle{Informativa Privacy}
\pnSubtitle{per i candidati al Recruitment}

\begin{document}

\pnMakeHeader

\section{Titolare del trattamento}

Ai sensi degli artt. 4 e 24 del Reg. UE 2016/679 il titolare del trattamento è individuabile nella figura del legale rappresentante dell'associazione PoliNetwork APS, il suo Presidente.

L'associazione non ha nominato un Responsabile della Protezione dei Dati (DPO) in quanto non ricorrono le condizioni di cui all'art. 37 del GDPR.

\textbf{Contatti per esercitare i diritti privacy:} \\
\textbf{Indirizzo email:} \href{mailto:privacy@polinetwork.org}{privacy@polinetwork.org} (in alternativa \href{mailto:direttivo@polinetwork.org}{direttivo@polinetwork.org}) \\
\textbf{PEC:} \href{mailto:polinetwork@pec.it}{polinetwork@pec.it}

\section{Tipologia di dati raccolti}

Tramite il modulo di candidatura "Join the Team!" raccogliamo e trattiamo i seguenti dati personali:
\begin{itemize}
    \item \textbf{Dati anagrafici e di contatto:} nome, cognome, indirizzo email istituzionale messo a disposizione dal Politecnico di Milano, username impostato sulla piattaforma Telegram, numero di telefono, username LinkedIn e preferenza sul canale di ricontatto.
    \item \textbf{Dati accademici:} anno frequentato, corso di studi e sede del Politecnico di Milano.
    \item \textbf{Dati curriculari e motivazionali:} precedenti o attuali esperienze in altre associazioni studentesche, aree di interesse per il contributo in associazione (IT, Design, Social Media, HR, ecc.), motivazioni della candidatura e link a portfolio o repository GitHub forniti volontariamente.
    \item \textbf{Dati statistici:} canale tramite il quale si è venuti a conoscenza del form.
\end{itemize}

\section{Finalità e base giuridica del trattamento}

I dati personali saranno trattati esclusivamente per le seguenti finalità:
\begin{enumerate}
    \item gestione delle candidature e del processo di selezione dei nuovi volontari;
    \item comunicazioni con i candidati relative al processo di selezione in corso o futuro;
    \item eventuale instaurazione del rapporto di collaborazione o ingresso in associazione.
\end{enumerate}
La base giuridica del trattamento è il consenso esplicito del candidato, espresso contestualmente all'invio del modulo di candidatura (art. 6, par. 1, lett. a del GDPR).

Il conferimento dei dati è facoltativo. Tuttavia, il mancato conferimento dei dati contrassegnati come obbligatori nel modulo renderà impossibile procedere alla valutazione della candidatura. Il conferimento dei dati facoltativi (ad es. link a portfolio o profilo LinkedIn) non ha conseguenze sulla procedura di selezione.

\section{Modalità del trattamento}

Il trattamento dei dati personali avviene mediante strumenti manuali, informatici e telematici. Le logiche applicate sono strettamente correlate alle finalità stesse e implementate in modo da garantire la massima sicurezza e riservatezza dei dati, prevenendone la perdita, gli usi illeciti e gli accessi non autorizzati.

\section{Conservazione dei dati}

I dati personali saranno conservati per il tempo strettamente necessario a conseguire le finalità per cui sono stati raccolti. Nello specifico:
\begin{itemize}
    \item i dati saranno conservati in forma attiva per l'intera durata del processo di selezione;
    \item per i candidati \textbf{non ammessi}: al termine della selezione, i dati essenziali relativi all'esito e all'andamento della candidatura potranno essere conservati per un periodo massimo di \textbf{due anni}, al solo fine di poter valutare eventuali candidature successive da parte della stessa persona e garantire la coerenza del processo selettivo nel tempo. Al termine di questo periodo i dati saranno cancellati o resi completamente anonimi;
    \item per i candidati \textbf{ammessi} che entrano a far parte dell'associazione: i dati confluiranno nel fascicolo del volontario e saranno trattati secondo l'informativa privacy specifica per i membri dell'associazione, che sarà fornita al momento dell'ingresso.
\end{itemize}

\section{Destinatari e accesso ai dati}

I dati personali non saranno comunicati a soggetti terzi esterni all'associazione né diffusi pubblicamente. I dati potranno essere accessibili esclusivamente ai membri dell'associazione PoliNetwork APS direttamente coinvolti nel processo di selezione (ad esempio membri dell'area HR e referenti delle aree tecniche di interesse del candidato), che tratteranno i dati nel pieno rispetto delle direttive interne di riservatezza.

Potranno avere accesso ai dati anche i soggetti che forniscono servizi per la gestione dei sistemi informatici del Titolare, i quali operano in qualità di responsabili del trattamento ai sensi dell'art. 28 del GDPR. In particolare, per la raccolta delle candidature viene utilizzata la piattaforma \textbf{Microsoft Forms}, fornita da Microsoft Corporation.

\textbf{Trasferimento dei dati verso Paesi terzi.} Microsoft Corporation può trattare dati personali negli Stati Uniti d'America. Tale trasferimento avviene nel rispetto del GDPR sulla base delle seguenti garanzie:
\begin{itemize}
    \item l'adesione di Microsoft al \textbf{Data Privacy Framework (DPF) UE-USA}, adottato dalla Commissione Europea con decisione di adeguatezza del 10 luglio 2023 (ai sensi dell'art. 45 del GDPR). L'elenco delle organizzazioni aderenti al DPF è consultabile al \href{https://www.dataprivacyframework.gov}{sito};
    \item la stipula di \textbf{Clausole Contrattuali Standard (SCC)} approvate dalla Commissione Europea (ai sensi dell'art. 46 del GDPR), che Microsoft incorpora nelle proprie condizioni contrattuali come misura di garanzia complementare.
\end{itemize}
Per ulteriori informazioni sulle modalità di trattamento dei dati da parte di Microsoft si rimanda alla relativa informativa privacy, consultabile all'indirizzo \href{https://privacy.microsoft.com}{privacy.microsoft.com}.

\section{Diritti degli interessati}

L'interessato potrà far valere i propri diritti come espressi dagli artt. 15-22 del Regolamento UE 2016/679, rivolgendosi al Titolare del Trattamento tramite i recapiti indicati nella Sezione 1 ("Titolare del trattamento") del presente documento.

Nello specifico, l'interessato ha il diritto, in qualunque momento, di:
\begin{itemize}
    \item \textbf{Diritto di accesso (art. 15):} ottenere la conferma dell'esistenza o meno di dati personali che lo riguardano, accedervi e riceverne copia;
    \item \textbf{Diritto di rettifica (art. 16):} richiedere l'aggiornamento, la rettifica di dati inesatti o l'integrazione di quelli incompleti;
    \item \textbf{Diritto alla cancellazione o "all'oblio" (art. 17):} ottenere la cancellazione dei propri dati qualora non siano più necessari per le finalità per le quali sono stati raccolti, o in caso di revoca del consenso;
    \item \textbf{Diritto di limitazione del trattamento (art. 18):} richiedere che i dati siano unicamente conservati (e non altrimenti trattati) in determinate ipotesi previste dalla legge;
    \item \textbf{Diritto alla portabilità dei dati (art. 20):} ricevere in un formato strutturato, di uso comune e leggibile da dispositivo automatico i dati personali forniti e, se tecnicamente fattibile, trasmetterli a un altro titolare senza impedimenti;
    \item \textbf{Diritto di opposizione (art. 21):} opporsi in tutto o in parte, per motivi connessi alla propria situazione particolare, al trattamento dei dati personali che lo riguardano.
\end{itemize}

Si informa inoltre che il processo di selezione \textbf{non è soggetto a decisioni basate unicamente sul trattamento automatizzato}, ivi compresa la profilazione, che producano effetti giuridici o che incidano significativamente sulla persona (art. 22 del GDPR). Le decisioni relative alle candidature sono sempre adottate da persone fisiche.

Fatto salvo ogni altro ricorso amministrativo e giurisdizionale, se il candidato ritiene che il trattamento dei dati violi quanto previsto dal GDPR, ha il diritto di proporre reclamo al Garante per la protezione dei dati personali ai sensi dell'art. 77 del GDPR (\href{https://www.garanteprivacy.it}{www.garanteprivacy.it}).

Infine, conformemente all'art. 7, paragrafo 3, e all'art. 6, paragrafo 1, lettera a) del GDPR, l'interessato ha il diritto di revocare il proprio consenso in qualsiasi momento, senza che ciò pregiudichi la liceità del trattamento basata sul consenso prestato prima della revoca.

\end{document}
